% ===== §3 =====
\section{The Living Present}
\label{sec:maintenance}

\noindent
This section takes the second dimension, the maintenance of a relation through time, and establishes that what maintenance depends on, the living present in which two people are actually with each other, cannot be occupied by any statement, because the act of stating already places the speaker at a distance from the moment stated. The result is that maintenance is a matter of presence rather than of statement, and the section shows this by way of the phenomenon, a review of the two accounts that would make the present sayable, the isolation of the core, and the claim.

\subsection*{The phenomenon}

A relation is not only brought into being; it is kept. Maintenance is the work of holding a bond alive through time, and what it holds is not a record of the past but a living present, the felt now in which two people are actually with each other. Language seems able to reach this present. We have the tense for it. I love you now. I am here. This is the present indicative, the grammatical form of the current moment, and it appears made to seize exactly what maintenance depends on.

It does not seize it. By the time \emph{I love you now} has been said, the now it names has slipped past. The present of feeling, the actual current in which the relation is alive, is always a step ahead of the signifier that would fix it, so that speech about the present is always speech about a present just departed. The word arrives late. What it catches is not the living now but the trace of a now already becoming past, and the living now itself, the thing maintenance most needs to hold, is precisely what runs out from under every attempt to state it.

\subsection*{Two accounts of the sayable present, and their remainder}

Two accounts promise that the present can, after all, be reached in speech, and each illuminates the limit by the exact manner of its failure.

The first is the account of the performative present, the deictic now. Linguistics observes that words like \emph{now}, \emph{here}, \emph{I}, and \emph{this} take their reference from the very act of utterance; the \emph{now} of \emph{I love you now} just is the moment of speaking, so it seems the word cannot lag the moment it names, since it names whatever moment it is spoken in. But this only relocates the lag. The deictic \emph{now} refers to the moment of utterance, yet the living present that maintenance needs is not the instant of an utterance; it is the felt current of being-with, which is not itself an act of speech and which the utterance interrupts in order to name. When I stop being with you in order to say that I am with you, the saying occupies the moment the deixis points to, and the being-with it would report has been set aside for the duration of the report. The deictic present is real, but it is the present of the speaking, not the present of the bond, and these are not the same present.

The second is the phenomenological account of the living present, which comes nearer the truth and for that reason marks the limit more precisely. Phenomenology describes the present of experience not as a knife-edge instant but as a thick now that carries the just-past in retention and the about-to-come in protention, so that the living present is always already a passage rather than a point. Merleau-Ponty makes perception itself this temporal thickness, a present that is never a bare instant but a field in motion \citep{merleauponty1962phenomenology}. This is exactly right about the present that maintenance inhabits, and it is precisely what defeats the statement. A statement fixes; it arrests a content and holds it still for predication. The living present, being a passage and not a point, cannot be arrested without being stopped, and stopped it is no longer the living present but a still taken from it. Phenomenology thus tells us why the present cannot be said: to say it is to fix it, and to fix a passage is to end the passing that made it alive.

\subsection*{The core}

The core language cannot reach here is the \emph{present} of the relation, its currentness, the being-now of the bond as against every account of how it was a moment ago. The relation lives in a present that speech cannot occupy, because the act of speaking already places the speaker at a small distance from the moment spoken of, and because the moment in question is a passage that fixing would end. One can narrate the relation, and narration has its place, but narration is always of the just-past, and the living edge on which the bond is actually sustained stays ahead of the narration, uncaptured.

\begin{claim}[The uncapturability of the present]
The living present in which a relation is sustained cannot be occupied by any statement. The deictic present belongs to the act of speaking and not to the being-with it would report; the thick present of experience is a passage that fixing would arrest and so end. What maintenance rests on stays ahead of every sentence that reaches for it.
\end{claim}

\begin{casebox}{the photograph of the moment}
Two people try to hold an evening by recording it, a photograph, a note written the same night, a message that says exactly how it feels now. Later the record is faithful and the moment is gone from it, present only as something the record points back to and does not contain. The living being-together that the evening was cannot be found in what was saved of it. This is not the failure of a poor record. The best record would do no better, because the present it would hold is a passage, and what is saved of a passage is always a still from which the passing has been removed.
\end{casebox}

\subsection*{Why the limit is generative here}

This is what makes maintenance a matter of presence rather than of statement. Because the living present cannot be said, it can only be inhabited, and the keeping of a relation is not the repeated production of true sentences about it but the repeated occupation of a present that no sentence holds. Consider again the counter-case, a bond whose current moment \emph{could} be fixed in words. Its maintenance would consist in keeping the words up to date, in issuing at each moment the true sentence for that moment, a maintenance of the record rather than of the bond; and a relation maintained by keeping its record current would be one in which presence had been replaced by reporting. The limit forbids this substitution. Because the living present cannot be caught in a sentence, maintenance cannot be discharged by sentences, and is thrown back onto presence itself, onto the being-there that no report replaces. The uncapturability that looked like a failure of the present tense is what keeps maintenance an inhabiting rather than a filing. It is also the temporal form of desire, which slides from signifier to signifier precisely because no signifier ever coincides with the present it reaches for, a movement the paper takes up in Part Two.
